\item (a) Considere que la definición del coeficiente de expansión volumétrica es
$$\beta = \frac{1}{V} \left. \frac{dV}{dT} \right|_{P = \text{constante}} = \frac{1}{V} \frac{\partial V}{\partial T}$$
Use la ecuación de estado de un gas ideal para demostrar que el coeficiente de expansión volumétrica de un gas ideal a presión constante está dado por $\beta = 1/T$, donde $T$ es la temperatura absoluta. (b) ¿Qué valor predice esta expresión para $\beta$ a 0 °C? Establezca cómo se compara este resultado con los valores experimentales para (c) helio y (d) aire en la tabla 6.1. \textit{Nota:} Estos valores son mayores que los coeficientes de expansión volumétrica para la mayoría de los líquidos y sólidos.
